
\slajdid{05-s037}
% element: 05-s037:e01
% element: 05-s037:e02
% element: 05-s037:e03
\begin{frame}[label=05-s037]{Neoznačeno množenje}
\fontsize{10}{12}\selectfont\setlength{\parskip}{0pt}\setlength{\abovedisplayskip}{2pt}\setlength{\belowdisplayskip}{2pt}

Neoznačeno\qquad $5\mathrm{x}3=101\mathrm{x}011=15=1111$
\begin{columns}[c,onlytextwidth]\column{.67\textwidth}\centering\begin{tikzpicture}[x=.55cm,y=.43cm,font=\fontsize{9}{10}\selectfont]
\filldraw[fill=DEplava,draw=white,line width=.4pt] (0,0) rectangle ++(1,-1);
\filldraw[fill=DEplava,draw=white,line width=.4pt] (1,0) rectangle ++(1,-1);
\filldraw[fill=DEplava,draw=white,line width=.4pt] (2,0) rectangle ++(1,-1);
\filldraw[fill=DEplava,draw=white,line width=.4pt] (3,0) rectangle ++(1,-1);
\node[text=white] at (3.5,-0.5) {$1$};
\filldraw[fill=DEplava,draw=white,line width=.4pt] (4,0) rectangle ++(1,-1);
\node[text=white] at (4.5,-0.5) {$0$};
\filldraw[fill=DEplava,draw=white,line width=.4pt] (5,0) rectangle ++(1,-1);
\node[text=white] at (5.5,-0.5) {$1$};
\filldraw[fill=DEplava,draw=white,line width=.4pt] (6,0) rectangle ++(1,-1);
\node[text=white] at (6.5,-0.5) {$x$};
\filldraw[fill=DEplava,draw=white,line width=.4pt] (7,0) rectangle ++(1,-1);
\node[text=white] at (7.5,-0.5) {$0$};
\filldraw[fill=DEplava,draw=white,line width=.4pt] (8,0) rectangle ++(1,-1);
\node[text=white] at (8.5,-0.5) {$1$};
\filldraw[fill=DEplava,draw=white,line width=.4pt] (9,0) rectangle ++(1,-1);
\node[text=white] at (9.5,-0.5) {$1$};
\filldraw[fill=DEplava!12,draw=white,line width=.4pt] (0,-1) rectangle ++(1,-1);
\filldraw[fill=DEplava!12,draw=white,line width=.4pt] (1,-1) rectangle ++(1,-1);
\filldraw[fill=DEplava!12,draw=white,line width=.4pt] (2,-1) rectangle ++(1,-1);
\filldraw[fill=DEplava!12,draw=white,line width=.4pt] (3,-1) rectangle ++(1,-1);
\node[text=DEtekst] at (3.5,-1.5) {$1$};
\filldraw[fill=DEplava!12,draw=white,line width=.4pt] (4,-1) rectangle ++(1,-1);
\node[text=DEtekst] at (4.5,-1.5) {$0$};
\filldraw[fill=DEplava!12,draw=white,line width=.4pt] (5,-1) rectangle ++(1,-1);
\node[text=DEtekst] at (5.5,-1.5) {$1$};
\filldraw[fill=DEplava!12,draw=white,line width=.4pt] (6,-1) rectangle ++(1,-1);
\filldraw[fill=DEplava!12,draw=white,line width=.4pt] (7,-1) rectangle ++(1,-1);
\filldraw[fill=DEplava!12,draw=white,line width=.4pt] (8,-1) rectangle ++(1,-1);
\filldraw[fill=DEplava!12,draw=white,line width=.4pt] (9,-1) rectangle ++(1,-1);
\filldraw[fill=DEplava!7,draw=white,line width=.4pt] (0,-2) rectangle ++(1,-1);
\node[text=DEtekst] at (0.5,-2.5) {$+$};
\filldraw[fill=DEplava!7,draw=white,line width=.4pt] (1,-2) rectangle ++(1,-1);
\filldraw[fill=DEplava!7,draw=white,line width=.4pt] (2,-2) rectangle ++(1,-1);
\node[text=DEtekst] at (2.5,-2.5) {$1$};
\filldraw[fill=DEplava!7,draw=white,line width=.4pt] (3,-2) rectangle ++(1,-1);
\node[text=DEtekst] at (3.5,-2.5) {$0$};
\filldraw[fill=DEplava!7,draw=white,line width=.4pt] (4,-2) rectangle ++(1,-1);
\node[text=DEtekst] at (4.5,-2.5) {$1$};
\filldraw[fill=DEplava!7,draw=white,line width=.4pt] (5,-2) rectangle ++(1,-1);
\filldraw[fill=DEplava!7,draw=white,line width=.4pt] (6,-2) rectangle ++(1,-1);
\filldraw[fill=DEplava!7,draw=white,line width=.4pt] (7,-2) rectangle ++(1,-1);
\filldraw[fill=DEplava!7,draw=white,line width=.4pt] (8,-2) rectangle ++(1,-1);
\filldraw[fill=DEplava!7,draw=white,line width=.4pt] (9,-2) rectangle ++(1,-1);
\filldraw[fill=DEplava!12,draw=white,line width=.4pt] (0,-3) rectangle ++(1,-1);
\node[text=DEtekst] at (0.5,-3.5) {$+$};
\filldraw[fill=DEplava!12,draw=white,line width=.4pt] (1,-3) rectangle ++(1,-1);
\node[text=DEtekst] at (1.5,-3.5) {$0$};
\filldraw[fill=DEplava!12,draw=white,line width=.4pt] (2,-3) rectangle ++(1,-1);
\node[text=DEtekst] at (2.5,-3.5) {$0$};
\filldraw[fill=DEplava!12,draw=white,line width=.4pt] (3,-3) rectangle ++(1,-1);
\node[text=DEtekst] at (3.5,-3.5) {$0$};
\filldraw[fill=DEplava!12,draw=white,line width=.4pt] (4,-3) rectangle ++(1,-1);
\filldraw[fill=DEplava!12,draw=white,line width=.4pt] (5,-3) rectangle ++(1,-1);
\filldraw[fill=DEplava!12,draw=white,line width=.4pt] (6,-3) rectangle ++(1,-1);
\filldraw[fill=DEplava!12,draw=white,line width=.4pt] (7,-3) rectangle ++(1,-1);
\filldraw[fill=DEplava!12,draw=white,line width=.4pt] (8,-3) rectangle ++(1,-1);
\filldraw[fill=DEplava!12,draw=white,line width=.4pt] (9,-3) rectangle ++(1,-1);
\filldraw[fill=DEplava!7,draw=white,line width=.4pt] (0,-4) rectangle ++(1,-1);
\filldraw[fill=DEplava!7,draw=white,line width=.4pt] (1,-4) rectangle ++(1,-1);
\node[text=DEtekst] at (1.5,-4.5) {$0$};
\filldraw[fill=DEplava!7,draw=white,line width=.4pt] (2,-4) rectangle ++(1,-1);
\node[text=DEtekst] at (2.5,-4.5) {$1$};
\filldraw[fill=DEplava!7,draw=white,line width=.4pt] (3,-4) rectangle ++(1,-1);
\node[text=DEtekst] at (3.5,-4.5) {$1$};
\filldraw[fill=DEplava!7,draw=white,line width=.4pt] (4,-4) rectangle ++(1,-1);
\node[text=DEtekst] at (4.5,-4.5) {$1$};
\filldraw[fill=DEplava!7,draw=white,line width=.4pt] (5,-4) rectangle ++(1,-1);
\node[text=DEtekst] at (5.5,-4.5) {$1$};
\filldraw[fill=DEplava!7,draw=white,line width=.4pt] (6,-4) rectangle ++(1,-1);
\filldraw[fill=DEplava!7,draw=white,line width=.4pt] (7,-4) rectangle ++(1,-1);
\filldraw[fill=DEplava!7,draw=white,line width=.4pt] (8,-4) rectangle ++(1,-1);
\filldraw[fill=DEplava!7,draw=white,line width=.4pt] (9,-4) rectangle ++(1,-1);
\draw[DEtekst] (1,-4)--(6,-4);
\draw[DEplava,->] (9.5,-1)--(9.5,-1.35)--(6.3,-1.35);
\draw[DEplava,->] (8.5,-1)--(8.5,-2.35)--(6.3,-2.35);
\draw[DEplava,->] (7.5,-1)--(7.5,-3.35)--(6.3,-3.35);
\end{tikzpicture}

\column{.30\textwidth}$\left\}\begin{gathered}\text{Parcijalni proizvodi}\\\text{logička I operacija}\\\text{bita operanada}\end{gathered}\right.$
\end{columns}\par\vspace{1mm}
\begin{columns}[c,onlytextwidth]\column{.67\textwidth}\centering\begin{tikzpicture}[x=.55cm,y=.43cm,font=\fontsize{9}{10}\selectfont]
\filldraw[fill=DEplava,draw=white,line width=.4pt] (0,0) rectangle ++(1,-1);
\filldraw[fill=DEplava,draw=white,line width=.4pt] (1,0) rectangle ++(1,-1);
\filldraw[fill=DEplava,draw=white,line width=.4pt] (2,0) rectangle ++(1,-1);
\filldraw[fill=DEplava,draw=white,line width=.4pt] (3,0) rectangle ++(1,-1);
\filldraw[fill=DEplava,draw=white,line width=.4pt] (4,0) rectangle ++(1,-1);
\node[text=white] at (4.5,-0.5) {$1$};
\filldraw[fill=DEplava,draw=white,line width=.4pt] (5,0) rectangle ++(1,-1);
\node[text=white] at (5.5,-0.5) {$0$};
\filldraw[fill=DEplava,draw=white,line width=.4pt] (6,0) rectangle ++(1,-1);
\node[text=white] at (6.5,-0.5) {$1$};
\filldraw[fill=DEplava,draw=white,line width=.4pt] (7,0) rectangle ++(1,-1);
\node[text=white] at (7.5,-0.5) {$x$};
\filldraw[fill=DEplava,draw=white,line width=.4pt] (8,0) rectangle ++(1,-1);
\node[text=white] at (8.5,-0.5) {$0$};
\filldraw[fill=DEplava,draw=white,line width=.4pt] (9,0) rectangle ++(1,-1);
\node[text=white] at (9.5,-0.5) {$1$};
\filldraw[fill=DEplava,draw=white,line width=.4pt] (10,0) rectangle ++(1,-1);
\node[text=white] at (10.5,-0.5) {$1$};
\filldraw[fill=DEplava!12,draw=white,line width=.4pt] (0,-1) rectangle ++(1,-1);
\filldraw[fill=DEplava!12,draw=white,line width=.4pt] (1,-1) rectangle ++(1,-1);
\node[text=DEcrvena] at (1.5,-1.5) {$0$};
\filldraw[fill=DEplava!12,draw=white,line width=.4pt] (2,-1) rectangle ++(1,-1);
\node[text=DEcrvena] at (2.5,-1.5) {$0$};
\filldraw[fill=DEplava!12,draw=white,line width=.4pt] (3,-1) rectangle ++(1,-1);
\node[text=DEcrvena] at (3.5,-1.5) {$0$};
\filldraw[fill=DEplava!12,draw=white,line width=.4pt] (4,-1) rectangle ++(1,-1);
\node[text=DEtekst] at (4.5,-1.5) {$1$};
\filldraw[fill=DEplava!12,draw=white,line width=.4pt] (5,-1) rectangle ++(1,-1);
\node[text=DEtekst] at (5.5,-1.5) {$0$};
\filldraw[fill=DEplava!12,draw=white,line width=.4pt] (6,-1) rectangle ++(1,-1);
\node[text=DEtekst] at (6.5,-1.5) {$1$};
\filldraw[fill=DEplava!12,draw=white,line width=.4pt] (7,-1) rectangle ++(1,-1);
\filldraw[fill=DEplava!12,draw=white,line width=.4pt] (8,-1) rectangle ++(1,-1);
\filldraw[fill=DEplava!12,draw=white,line width=.4pt] (9,-1) rectangle ++(1,-1);
\filldraw[fill=DEplava!12,draw=white,line width=.4pt] (10,-1) rectangle ++(1,-1);
\filldraw[fill=DEplava!7,draw=white,line width=.4pt] (0,-2) rectangle ++(1,-1);
\node[text=DEtekst] at (0.5,-2.5) {$+$};
\filldraw[fill=DEplava!7,draw=white,line width=.4pt] (1,-2) rectangle ++(1,-1);
\node[text=DEcrvena] at (1.5,-2.5) {$0$};
\filldraw[fill=DEplava!7,draw=white,line width=.4pt] (2,-2) rectangle ++(1,-1);
\node[text=DEcrvena] at (2.5,-2.5) {$0$};
\filldraw[fill=DEplava!7,draw=white,line width=.4pt] (3,-2) rectangle ++(1,-1);
\node[text=DEtekst] at (3.5,-2.5) {$1$};
\filldraw[fill=DEplava!7,draw=white,line width=.4pt] (4,-2) rectangle ++(1,-1);
\node[text=DEtekst] at (4.5,-2.5) {$0$};
\filldraw[fill=DEplava!7,draw=white,line width=.4pt] (5,-2) rectangle ++(1,-1);
\node[text=DEtekst] at (5.5,-2.5) {$1$};
\filldraw[fill=DEplava!7,draw=white,line width=.4pt] (6,-2) rectangle ++(1,-1);
\node[text=DEisticanje] at (6.5,-2.5) {$0$};
\filldraw[fill=DEplava!7,draw=white,line width=.4pt] (7,-2) rectangle ++(1,-1);
\filldraw[fill=DEplava!7,draw=white,line width=.4pt] (8,-2) rectangle ++(1,-1);
\filldraw[fill=DEplava!7,draw=white,line width=.4pt] (9,-2) rectangle ++(1,-1);
\filldraw[fill=DEplava!7,draw=white,line width=.4pt] (10,-2) rectangle ++(1,-1);
\filldraw[fill=DEplava!12,draw=white,line width=.4pt] (0,-3) rectangle ++(1,-1);
\node[text=DEtekst] at (0.5,-3.5) {$+$};
\filldraw[fill=DEplava!12,draw=white,line width=.4pt] (1,-3) rectangle ++(1,-1);
\node[text=DEcrvena] at (1.5,-3.5) {$0$};
\filldraw[fill=DEplava!12,draw=white,line width=.4pt] (2,-3) rectangle ++(1,-1);
\node[text=DEtekst] at (2.5,-3.5) {$0$};
\filldraw[fill=DEplava!12,draw=white,line width=.4pt] (3,-3) rectangle ++(1,-1);
\node[text=DEtekst] at (3.5,-3.5) {$0$};
\filldraw[fill=DEplava!12,draw=white,line width=.4pt] (4,-3) rectangle ++(1,-1);
\node[text=DEtekst] at (4.5,-3.5) {$0$};
\filldraw[fill=DEplava!12,draw=white,line width=.4pt] (5,-3) rectangle ++(1,-1);
\node[text=DEisticanje] at (5.5,-3.5) {$0$};
\filldraw[fill=DEplava!12,draw=white,line width=.4pt] (6,-3) rectangle ++(1,-1);
\node[text=DEisticanje] at (6.5,-3.5) {$0$};
\filldraw[fill=DEplava!12,draw=white,line width=.4pt] (7,-3) rectangle ++(1,-1);
\filldraw[fill=DEplava!12,draw=white,line width=.4pt] (8,-3) rectangle ++(1,-1);
\filldraw[fill=DEplava!12,draw=white,line width=.4pt] (9,-3) rectangle ++(1,-1);
\filldraw[fill=DEplava!12,draw=white,line width=.4pt] (10,-3) rectangle ++(1,-1);
\filldraw[fill=DEplava!7,draw=white,line width=.4pt] (0,-4) rectangle ++(1,-1);
\filldraw[fill=DEplava!7,draw=white,line width=.4pt] (1,-4) rectangle ++(1,-1);
\node[text=DEtekst] at (1.5,-4.5) {$0$};
\filldraw[fill=DEplava!7,draw=white,line width=.4pt] (2,-4) rectangle ++(1,-1);
\node[text=DEtekst] at (2.5,-4.5) {$0$};
\filldraw[fill=DEplava!7,draw=white,line width=.4pt] (3,-4) rectangle ++(1,-1);
\node[text=DEtekst] at (3.5,-4.5) {$1$};
\filldraw[fill=DEplava!7,draw=white,line width=.4pt] (4,-4) rectangle ++(1,-1);
\node[text=DEtekst] at (4.5,-4.5) {$1$};
\filldraw[fill=DEplava!7,draw=white,line width=.4pt] (5,-4) rectangle ++(1,-1);
\node[text=DEtekst] at (5.5,-4.5) {$1$};
\filldraw[fill=DEplava!7,draw=white,line width=.4pt] (6,-4) rectangle ++(1,-1);
\node[text=DEtekst] at (6.5,-4.5) {$1$};
\filldraw[fill=DEplava!7,draw=white,line width=.4pt] (7,-4) rectangle ++(1,-1);
\filldraw[fill=DEplava!7,draw=white,line width=.4pt] (8,-4) rectangle ++(1,-1);
\filldraw[fill=DEplava!7,draw=white,line width=.4pt] (9,-4) rectangle ++(1,-1);
\filldraw[fill=DEplava!7,draw=white,line width=.4pt] (10,-4) rectangle ++(1,-1);
\draw[DEtekst] (1,-4)--(7,-4);
\end{tikzpicture}

\column{.30\textwidth}\textcolor{DEisticanje}{0} – ASL, x2\\\textcolor{DEcrvena}{0} - Proširenje opsega
\end{columns}
\[(2^n-1)(2^n-1)=(\textcolor{DEcrvena}{2^{2n}-1})-(2^{n+1}-1)+1\]
\begin{center}Treba nam $2n$ mesta za smeštanje rezultata\end{center}
\end{frame}
